% Generated draft: requires mathematical and source review.
\documentclass[11pt]{article}
\usepackage[margin=25mm]{geometry}
\usepackage{fontspec}
\setmainfont{Latin Modern Roman}
\usepackage{amsmath,amssymb,mathtools,unicode-math}
\setmathfont{Latin Modern Math}
\usepackage{enumitem,needspace,xurl,hyperref}
\setlength{\parindent}{0pt}
\setlength{\parskip}{5pt}
\setlength{\emergencystretch}{3em}
\newcommand{\N}{\mathbb{N}}
\newcommand{\Z}{\mathbb{Z}}
\newcommand{\Q}{\mathbb{Q}}
\newcommand{\R}{\mathbb{R}}
\newcommand{\C}{\mathbb{C}}
\newcommand{\F}{\mathbb{F}}
\newcommand{\A}{\mathbb{A}}
\newcommand{\PP}{\mathbb{P}}
\newcommand{\E}{\mathbb{E}}
\newcommand{\eps}{\varepsilon}
\newcommand{\dd}{\,\mathrm{d}}
\begin{document}
\section*{UnsolvedMath ID: 20001916}
\textbf{Source problem code:} catalogue entry
\subsection{Definitions and mathematical statement}
Let \(M\) be a smooth closed connected orientable surface of genus at least one, with a smooth Riemannian metric \(g\). Question 8.5.2 asks whether there always exists a complete unit-speed geodesic \(\gamma:\R\to M\) such that
\[
\overline{\gamma(\R)}=M.
\]
Density here concerns the projected image in the surface, not the lifted orbit \((\gamma(t),\dot\gamma(t))\) in \(SM\). A forward-dense ray, \(\overline{\gamma([0,\infty))}=M\), gives a stronger affirmative example.

The supplied research aid reports forward-dense geodesics for every cyclic torus metric
\[
M=\R^2/(a\Z\times b\Z),\quad
g=\dd x^2+f(x)^2\dd y^2,\qquad a,b>0,
\]
where \(f\) is smooth, positive and \(a\)-periodic. The conserved Clairaut momentum \(J=f(x)^2\dot y\) is a nonconstant continuous function on \(SM\), so every orbit lies in one of its proper level sets and no orbit is dense in all of \(SM\). This obstruction is consistent with a geodesic dense in the base torus.

More generally, the aid reports the sufficient measure-theoretic criterion: if \(\nu\) is an ergodic geodesic-flow-invariant probability measure on \(SM\) and its base marginal \(\pi_*\nu\), for \(\pi:SM\to M\), has full support, then \(\nu\)-almost every orbit projects forward densely in \(M\). Full support of \(\pi_*\nu\) means positive measure for every nonempty open subset of \(M\); it does not require full support of \(\nu\) on \(SM\). The aid says that, combined with published results, this answers the question for all positive-genus surface metrics without conjugate points. Those published results are not identified, and the unrestricted question is not claimed solved.

The extract also begins section 8.6 with a separate theorem attributed to Gaidukov. For every Riemannian metric on an oriented closed surface of genus at least one, every nontrivial free homotopy class \(\Gamma\), and every \(p\in M\), it states the existence of a periodic geodesic \(\eta:\R\to M\) in \(\Gamma\) and a geodesic half-line \(\beta:[0,\infty)\to M\), \(\beta(0)=p\), with synchronized parametrizations satisfying
\[
d_g(\beta(t),\eta(t))\longrightarrow0\qquad(t\to\infty).
\]
Here ``ray'' is not taken to mean globally distance-minimizing in the compact base surface. The source attributes derivation from Morse and Hedlund to [Bia-Pol1986]. This asymptotic-to-a-periodic-geodesic statement is adjacent background; it does not imply a dense geodesic in \(M\).
\subsection{Short English statement}
Does every positive-genus surface metric have a geodesic dense in the surface? The supplied aid gives cyclic torus examples and an invariant-measure criterion, while distinguishing base density from a dense orbit in the full unit tangent bundle. The appended Gaidukov theorem is separate background.
\subsection{Sources}
\begin{itemize}
\item Keith Burns and Vladimir Matveev, \emph{Open problems and questions about geodesics}, Question 8.5.2: \url{https://aimath.org/pastworkshops/geodesicsproblems.pdf}.
\item Adjacent section 8.6 cites [Gaidukov1966], [Bia-Pol1986], [Morse1924] and [Hedlund1932]. Full entries are absent. The research aid's published no-conjugate-points results are not identified.
\item Local registry \texttt{data/unsolvedmath/problems.json}, permanent ID 20001916, code \texttt{catalogue entry}; canonical \texttt{aim-geometry-notes.json}, \texttt{notes[253]}, source item 8.5; updated 14 August 2026. The supplied verification note reports earlier source-PDF and adjacent-record inspection.
\item Research attempt 1: partial result, problem status at run open, rigor label proved, novelty confidence low; current registry status partially solved. The aid is machine generated, not peer reviewed, and requires independent verification before citation.
\end{itemize}
\end{document}
